\section{Componente di retrieval}\label{sec:retrieval}

A ogni domanda, il componente Retrieving deve individuare, tra i circa 2.500 frammenti indicizzati nel prototipo
(Sezione~\ref{sec:chunking}), i pochi davvero pertinenti da passare al modello generativo. Non è un singolo passo,
ma un imbuto a più stadi, via via più costosi e più precisi (Figura~\ref{fig:imbuto-retrieval}): prima la domanda
viene analizzata e, se necessario, riformulata; poi una ricerca ibrida seleziona un primo insieme ampio di
candidati; un modello di riordinamento li valuta uno a uno contro la domanda; infine i frammenti vengono fusi con
il loro contesto quando serve e filtrati per pertinenza. Questa sezione ne descrive l'adattamento concreto al dominio universitario.

\begin{figure}[H]
  \centering
  \begin{tikzpicture}[
      stg/.style={draw, thick, rounded corners=3pt, fill=blue!7, align=center, text width=2.3cm,
                  minimum height=1.1cm, font=\scriptsize\bfseries},
      sub/.style={font=\scriptsize, align=center, text=black!70, text width=2.5cm, anchor=north},
      arr/.style={-{Stealth[length=2mm]}, thick, black!70},
    ]
    \node[stg] (q) at (0,0) {Analisi della domanda};
    \node[stg] (hs) at (3.3,0) {Ricerca ibrida};
    \node[stg] (rr) at (6.6,0) {Riordinamento};
    \node[stg] (am) at (9.9,0) {Fusione e soglia};
    \foreach \a/\b in {q/hs, hs/rr, rr/am} \draw[arr] (\a) -- (\b);
    \node[sub] at (0,-0.9) {condensazione, scomposizione, corso};
    \node[sub] at (3.3,-0.9) {denso + sparso, RRF, $\sim$20 candidati};
    \node[sub] at (6.6,-0.9) {cross-encoder, priorità al corso};
    \node[sub] at (9.9,-0.9) {sezioni intere, soglia, i più pertinenti};
  \end{tikzpicture}
  \caption{I quattro stadi del recupero, ciascuno descritto in questa sezione. Un quinto meccanismo, HyDE,
  interviene solo quando il corso è noto e la ricerca ibrida non restituisce
  nulla.}
  \label{fig:imbuto-retrieval}
\end{figure}

\subsection{Analisi della domanda}\label{sec:analisi-domanda}

Prima di essere cercata, la domanda viene normalizzata in tre passi indipendenti. Il primo, la
\emph{condensazione}, riformula una domanda di follow-up in una domanda autonoma alla luce della cronologia della
conversazione, così che il recupero non debba mai cercare un
riferimento implicito. Il secondo, la \emph{scomposizione}, divide una domanda che unisce più richieste distinte in
altrettante domande autonome, cercate e poi fuse separatamente, perché il riordinamento valuta la pertinenza di un
frammento rispetto alla domanda intera, e un frammento pertinente a una sola delle due richieste rischia di uscire
dai primi posti quando l'altra richiesta, nello stesso testo, pesa quanto lei. Entrambe le tecniche sono discusse
nella Sezione~\ref{sec:rag}; nella pratica di questo sistema, la scomposizione richiede un'attenzione che la sola
teoria non rende evidente: quando una richiesta della domanda originale dipende implicitamente da quella
precedente, la domanda scomposta
deve ripetere esplicitamente a quale contesto fa riferimento: senza quel riferimento esplicito, la domanda potrebbe risultare
ambigua quanto quella originale. La stessa chiamata riconosce anche i messaggi che non sono domande
sull'università, come i saluti, che non vengono cercati e ricevono direttamente una risposta.

Il terzo passo, specifico di questo dominio, individua il corso di laurea nominato esplicitamente nella domanda.

Ad esempio, su una cronologia in cui lo studente ha già chiesto ``Entro quando scade il pagamento della
prima rata per gli studenti iscritti a una laurea triennale?'', la nuova domanda ``E per chi è iscritto alla magistrale, e quanto costa la mora se pago in ritardo?'' --- un follow-up che è anche una domanda composta --- viene
prima condensata dal modello linguistico in produzione in ``Ed entro quando scade il
pagamento della prima rata per gli studenti di laurea magistrale, e quanto costa la mora se pago
in ritardo?'', e poi scomposta nelle due domande autonome
\begin{quote}
  \small
  Entro quando scade il pagamento della prima rata per gli studenti di laurea magistrale?\\
  Quanto costa la mora se pago in ritardo la prima rata della laurea magistrale?
\end{quote}
cercate e fuse separatamente nello stadio successivo.

\subsection{Ricerca ibrida}\label{sec:ricerca-ibrida}

Alla domanda, già normalizzata (Sezione~\ref{sec:analisi-domanda}), si applica lo stesso modello di embedding usato
in indicizzazione (Sezione~\ref{sec:embedding}), e Qdrant esegue due ricerche in parallelo, una per la
rappresentazione densa e una per quella sparsa, fondendole con Reciprocal Rank Fusion (Sezione~\ref{sec:rag}). Quando il corso di laurea dello studente è noto, un filtro restringe la ricerca ai soli
frammenti di ateneo e a quelli di quel corso --- gli altri corsi, pur presenti nell'indice, non vengono nemmeno
candidati in questa fase. È una scelta deliberata, non solo un'ottimizzazione: più corsi pubblicano pagine quasi
omonime su argomenti analoghi, con contenuto diverso (ciascuno ha una propria \emph{Prova Pratica Valutativa}, una
propria pagina \emph{Tasse ed esenzioni}), e lasciarle tutte in gara rischierebbe di citare, nella stessa risposta,
requisiti o importi di corsi diversi. Quando il corso non è ancora noto, la ricerca resta invece ristretta alle sole
pagine di ateneo, finché la domanda --- o la pagina trovata --- non permette di risolverlo.
Il risultato di questa fase è un insieme di circa venti candidati, troppi per essere affidabili singolarmente ma
abbastanza pochi da poter essere valutati uno a uno nello stadio successivo.

Per la seconda delle due domande scomposte nell'esempio della Sezione~\ref{sec:analisi-domanda} --- ``Quanto costa
la mora se pago in ritardo la prima rata della laurea magistrale?'' --- ristretta al corso di Psicologia
Scolastica e di Comunità, la ricerca ibrida restituisce effettivamente venti candidati: il punteggio di fusione
decresce da 0,700 del primo a 0,159 dell'ottavo, un calo che mostra perché affidarsi al solo primo risultato della
ricerca ibrida sarebbe rischioso, ma anche perché valutare uno a uno con un modello più accurato solo una ventina di
candidati, invece di migliaia, resta praticabile a ogni domanda.

\subsection{Riordinamento e corso di laurea}\label{sec:corso-priorita}

I candidati della ricerca ibrida vengono valutati da un cross-encoder, \texttt{bge-reranker-v2-m3}
(Sezione~\ref{sec:rag}), che riceve la domanda e ciascun candidato insieme e ne stima la pertinenza reciproca --- più
accurato della sola similarità vettoriale, perché può cogliere corrispondenze specifiche tra i due testi che la
ricerca ibrida, confrontando rappresentazioni calcolate separatamente, non vede. Il candidato non è giudicato sul
solo testo del frammento, ma sul percorso delle sue intestazioni (Sezione~\ref{sec:chunking}) seguito dal testo,
esattamente come già avviene per l'embedding (Sezione~\ref{sec:embedding}): un frammento che è solo una tabella o un
elenco puntato, senza contesto, è spesso incomprensibile al cross-encoder da solo --- per la domanda ``Sotto quale
ISEE si ha l'esenzione?'', il punteggio del frammento ``l'indicatore \dots{} è pari o inferiore a 27.000 €'', che da
solo non dice a cosa si riferisca quella soglia, passa da 0,733 a 0,860 quando gli viene anteposta l'intestazione
della sua sezione. Se il corso di laurea è noto, prima
del riordinamento il suo nome viene rimosso dal testo della domanda: la ricerca è già ristretta a quel corso
(Sezione~\ref{sec:ricerca-ibrida}), e lasciarlo nella domanda farebbe sì che il cross-encoder, un modello sensibile
alla sovrapposizione lessicale, premiasse ogni frammento che nomina quel corso indipendentemente dal contenuto. Un
secondo meccanismo, un bonus fisso di 0,1 al punteggio dei frammenti del corso noto, interviene dopo il riordinamento: a
quel punto, con la ricerca già ristretta al solo corso noto e alle pagine di ateneo, il bonus non serve più a
tenere testa ad altri corsi --- non ce ne sono --- ma a preferire, quando entrambe compaiono tra i candidati, la
versione specifica del corso di un argomento rispetto alla versione generica di ateneo sullo stesso argomento.

Resta il caso in cui il recupero, ristretto al corso noto e alle pagine di ateneo, non trovi nulla sopra la soglia
nemmeno dopo un secondo tentativo con HyDE: può darsi che la domanda riguardi un altro
corso. In questo caso il sistema non prova a cercare da solo tra gli altri
corsi, ma lo chiede esplicitamente allo studente, offrendogli l'elenco dei corsi indicizzati; se lo studente ne
indica uno diverso, la ricerca riparte ristretta al nuovo corso.

Allo studente viene chiesto di indicare il proprio corso anche quando non è ancora noto affatto, ma per una ragione
diversa dalla precedente. In questo caso la ricerca usata per rispondere resta sempre ristretta alle sole pagine di
ateneo (Sezione~\ref{sec:ricerca-ibrida}): non è mai lei a cercare tra i corsi. Se non trova nulla sopra soglia,
allo studente viene chiesto il corso direttamente, come nel caso precedente. Se invece trova qualcosa, il sistema
lancia un'ulteriore ricerca, questa volta senza alcuna restrizione --- su tutto il corpus, corsi compresi --- usata
solo per un confronto e mai per rispondere: se tra i suoi risultati una pagina di un corso specifico ottiene un
punteggio più alto del miglior risultato trovato tra le sole pagine di ateneo, è segno che la domanda dipende da un
corso più di quanto le pagine generali possano dire, e anche in questo caso il sistema preferisce chiedere allo
studente il proprio corso piuttosto che rispondere con un contenuto di ateneo che potrebbe non riguardarlo.

\subsection{Fusione gerarchica e soglia di accettazione}

I frammenti riordinati attraversano infine l'\emph{AutoMergingRetriever} (Sezione~\ref{sec:rag}): quando abbastanza
fratelli della stessa sezione compaiono tra i candidati, vengono sostituiti dalla sezione intera che li contiene
(Sezione~\ref{sec:chunking}), così che il modello generativo riceva un contesto completo invece di frammenti
isolati e lessicalmente vicini tra loro. Solo a questo punto, sui nodi --- frammenti o sezioni fuse --- così
ottenuti, viene applicata una soglia minima di punteggio, 0,3 sui punteggi del riordinamento normalizzati tra 0 e
1: chi non la supera viene scartato, perché un contesto
debolmente pertinente, passato comunque al modello, aumenta il rischio di una risposta costruita su un frammento
sbagliato ma presente, invece di un'astensione corretta --- coerente con il vincolo
di precisione fattuale che guida l'intero sistema (Capitolo~\ref{cap:architettura}). Se nessun nodo raggiunge quella
soglia, un meccanismo di ripiego ne accetta comunque alcuni con una soglia più permissiva, 0,12, piuttosto che
restituire un contesto vuoto: un contesto vuoto forza il rifiuto anche nei casi in cui un frammento scartato dalla
soglia principale fosse in realtà quello giusto, solo formulato in un lessico che il riordinamento premia meno di
quanto meriterebbe. Questo non vuol dire accettare il rischio di una risposta sbagliata al posto di un'astensione:
il modello generativo resta comunque istruito a dichiarare esplicitamente quando il contesto ricevuto non contiene
la risposta, invece di inventarla, ed è questo secondo controllo --- più accurato
perché legge il testo del frammento invece del solo punteggio --- a decidere se un candidato del ripiego risponde
davvero o va scartato. Il ripiego allarga quindi l'insieme dei candidati sottoposti a quel controllo, senza
sostituirlo: la soglia più permissiva decide solo cosa il modello generativo può vedere, non cosa deve usare. Dei
nodi che superano la soglia, solo i più pertinenti vengono infine passati al modello generativo, in ordine di
punteggio decrescente.

\subsection{Recupero di riserva con HyDE}\label{sec:hyde-retrieval}

Quando il corso di laurea è noto ma la ricerca ibrida non restituisce alcun candidato sopra soglia, la domanda viene
cercata una seconda volta con HyDE (Sezione~\ref{sec:rag}): un paragrafo ipotetico, scritto dal modello linguistico
nel linguaggio amministrativo che una vera pagina d'ateneo userebbe, viene accodato alla domanda originale prima di
calcolarne l'embedding. Il caso tipico che questo risolve è un disallineamento di vocabolario tra lo studente e la
pagina che gli risponderebbe: una domanda sul ``test d'ingresso'' cercata contro una pagina che parla di ``prova di
ammissione'' ottiene comunque un punteggio alto nella ricerca densa, che coglie la vicinanza di significato, ma un
punteggio quasi nullo dal cross-encoder, contro lo 0,99 ottenuto dalla stessa pagina quando la query usa il
suo stesso lessico: resta sotto soglia, nonostante
il frammento giusto fosse stato trovato. Il paragrafo ipotetico, scritto con il lessico amministrativo atteso,
porta nella query proprio quella terminologia, mentre la domanda originale resta anch'essa nella query cercata, così
che il riordinamento continui a valutare la pertinenza rispetto a quanto è stato effettivamente chiesto. Al modello
non viene chiesto di essere accurato, solo verosimile nel linguaggio: le date che scrive non vengono mai mostrate
allo studente e servono solo a orientare la ricerca verso il lessico giusto.

La Figura~\ref{fig:hyde-esempio} ripercorre l'intero meccanismo su un caso reale: la domanda ``Quando si tiene il
test d'ingresso?'', il paragrafo ipotetico generato dal modello in produzione per quella domanda, e un frammento
realmente indicizzato --- la pagina del corso di Infermieristica sulla \emph{prova di ammissione} --- contro cui
entrambe le query vengono valutate dal cross-encoder. Senza HyDE, la domanda originale ottiene un punteggio di
0,03, ben sotto la soglia di accettazione; la query completata con il paragrafo ipotetico ottiene invece 0,94,
abbondantemente sopra soglia, pur non avendo il modello azzeccato il termine ``ammissione'' usato dalla pagina: la
sola maggiore ricchezza lessicale e informativa del paragrafo, nel linguaggio amministrativo atteso, è bastata a
farlo riconoscere come pertinente.

\begin{figure}[H]
  \centering
  \begin{tikzpicture}[
      box/.style={draw, thick, rounded corners=3pt, text width=6.2cm, font=\scriptsize, inner sep=5pt},
      lbl/.style={font=\tiny, align=center, text width=2.6cm},
      arr/.style={-{Stealth[length=2mm]}, thick},
    ]
    \node[box, fill=blue!7, align=center] (q) at (0,3.6)
      {\textbf{Domanda originale}\\ ``Quando si tiene il test d'ingresso?''};
    \node[box, fill=yellow!12, align=left] (hyde) at (0,1.15)
      {\textbf{Paragrafo ipotetico generato dal modello (HyDE)}\\[2pt]
       \itshape ``Il test d'ingresso \dots{} è previsto annualmente nel mese di luglio, con le sessioni che si
       svolgono dal 15 luglio al 31 luglio. Per partecipare, gli studenti devono effettuare la prenotazione
       online \dots{} entro il 10 luglio.''};
    \node[box, fill=blue!7, align=center] (qf) at (0,-1.5)
      {\textbf{Query con HyDE}\\ domanda originale + paragrafo ipotetico};
    \node[box, fill=green!6, align=left, text width=4.6cm] (page) at (9.3,1.1)
      {\textbf{Frammento indicizzato}\\ (corso di Infermieristica)\\[2pt]
       ``\dots{} la prova di ammissione per l'a.a.\ 2026/27 si svolgerà il giorno 16 settembre 2026 \dots{}''};

    \draw[arr, black!70] (q) -- (hyde);
    \draw[arr, black!70] (hyde) -- (qf);

    \draw[arr, red!60!black, dashed] (q.east) -- (page.north west)
      node[lbl, midway, above, sloped] {senza HyDE: 0,03\\ sotto soglia (0,3)};
    \draw[arr, green!45!black] (qf.east) -- (page.south west)
      node[lbl, midway, below, sloped] {con HyDE: 0,94\\ sopra soglia};
  \end{tikzpicture}
  \caption{Esempio reale del recupero di riserva con HyDE: la domanda originale, cercata da sola contro il
  frammento corretto, ottiene dal cross-encoder un punteggio troppo basso per superare la soglia; la stessa domanda,
  completata con il paragrafo ipotetico generato dal modello, lo supera ampiamente. Punteggi calcolati eseguendo dal
  vivo il reranker e il modello di HyDE in produzione contro un frammento reale dell'indice.}
  \label{fig:hyde-esempio}
\end{figure}
